\documentclass[10pt,a4paper]{article}
\usepackage[T1]{fontenc}
\usepackage[utf8]{inputenc}
\usepackage{lmodern,microtype,geometry,amsmath,amssymb,booktabs,array,tabularx,xcolor,enumitem,fancyhdr}
\definecolor{ink}{HTML}{17324D}
\definecolor{teal}{HTML}{127C80}
\usepackage[colorlinks=true,urlcolor=teal,linkcolor=teal]{hyperref}
\geometry{left=22mm,right=22mm,top=23mm,bottom=22mm,headheight=14pt}
\pagestyle{fancy}\fancyhf{}
\fancyhead[L]{\small\color{ink} DSC / Inference workflows}
\fancyhead[R]{\small 15 October 2026}
\fancyfoot[L]{\scriptsize Algorithm note; no new estimation results}
\fancyfoot[R]{\small\thepage\ / 3}
\setlength{\parindent}{0pt}\setlength{\parskip}{6pt}
\setlength{\emergencystretch}{2em}
\setlist[enumerate]{leftmargin=*,itemsep=3pt,topsep=3pt}
\renewcommand{\arraystretch}{1.23}
\newcolumntype{Y}{>{\raggedright\arraybackslash}X}
\newcommand{\Normal}{\mathcal N}
\newcommand{\IG}{\operatorname{IG}}
\newcommand{\IW}{\operatorname{IW}}
\newcommand{\code}[1]{\texttt{#1}}
\hypersetup{pdftitle={DSC: Estimation, smoothing and prior parameter sensitivity},pdfauthor={DSC research}}
\begin{document}
{\color{ink}\LARGE\bfseries Three ways to run the DSC model\par}
{\large Estimate parameters, smooth states, or hold parameter weights at the prior\par}
\medskip
The new prior option makes sense as a \textbf{prior-sensitivity exercise}: it shows
the correlation paths supported by the data when parameter sets retain their
prior weights. It answers a different question from posterior estimation.

\begin{tabularx}{\linewidth}{@{}>{\raggedright\arraybackslash}p{32mm}YY@{}}\toprule
Mode & Parameters & State paths \\\midrule
Estimation + smoothing & Learn parameters from returns; by default fix them at posterior means for stage 2. & Smooth using all available modeled returns. \\
Estimation only & Learn and save posterior parameter draws. & Sample states internally; skip the separate smoother and state figures. \\
Prior + smoothing & Draw parameter sets from the calibrated prior; keep them fixed. & Run separate conditional smoothers for each set; mix them with equal weights. \\\bottomrule
\end{tabularx}

\subsection*{What the model separates}
For weekly returns $y_t$, the latent states are the changing mean $B_t$,
log variances $h_t$, and correlation coordinates $r_t$. The fixed parameters
$\theta=(V,\sigma_h^2,\sigma_r^2)$ control how quickly these states move:
\begin{align*}
y_t &\sim\Normal(B_t,D_tP_tD_t), &D_t&=\operatorname{diag}\{e^{h_{jt}/2}\},\\
B_t-B_{t-1}&\sim\Normal(0,V), &
h_{jt}-h_{j,t-1}&\sim\Normal(0,\sigma_{h,j}^2),\\
r_{kt}-r_{k,t-1}&\sim\Normal(0,\sigma_{r,k}^2), &
r_t&=\operatorname{vecl}(\log P_t).
\end{align*}
$P_t$ is the actual correlation matrix; $r_t$ is an unrestricted numerical
coordinate. Smoothing uses future observations relative to a historical week,
so old paths may revise when the sample grows.

\subsection*{1. Current default: estimation, then conditional smoothing}
\begin{enumerate}
\item Reserve the initial 104 weeks for initial-prior calibration and exclude
them from the likelihood. Use modeled returns through the estimation cutoff.
\item Run the joint Metropolis-within-Gibbs sampler. A sweep draws $B$ by forward filtering and
backward simulation; updates $V$ from its inverse-Wishart conditional; updates
each $r$ path by elliptical slice sampling; then proposes $\sigma_r^2$ from
inverse-gamma distributions with a Metropolis correction for the initial-state
density. It updates $h$ and $\sigma_h^2$ in the same way.
\item Discard warm-up, retain draws, check diagnostics, and save the parameter
snapshot. The default mean-smoothing option takes posterior means of $\theta$.
\item Freeze those means and run state-only chains over the full modeled
history, targeting $p(B,h,r\mid y,\widehat\theta)$. No variance parameters are
updated in this stage.
\end{enumerate}
\textbf{Interpretation.} The default stage-2 bands contain conditional state
uncertainty; they exclude parameter uncertainty. Each of the two stages has
its own \code{MaxHours} limit. Diagnostics, not the stopping flag, determine
whether the saved draws are adequate.

\clearpage
{\color{ink}\Large\bfseries Estimation only, reuse, and calibration\par}
\subsection*{2. Estimate and save parameters only}
Set \code{ParameterSource="posterior"}, \code{EstimateParameters=true}, and
\code{SmoothStates=false}. Run the same joint sampler described on page 1,
check its diagnostics, and save the parameter snapshot. Skip the second
conditional smoother and the state figures.

\textbf{Why are states still sampled?} The conditional updates of $V$,
$\sigma_h^2$ and $\sigma_r^2$ use changes in $B$, $h$ and $r$. Estimation only
therefore means ``no separate smoothing/output stage,'' not ``no latent states
inside estimation.''

\begin{quote}\small\ttfamily
result = run\_weekly\_model\_report( ...\\
\hspace*{3mm}ParameterSource="posterior", EstimateParameters=true, ...\\
\hspace*{3mm}SmoothStates=false, WarmupIterations=100, ...\\
\hspace*{3mm}RetainedDraws=100, NumChains=4, ParallelChains=true, ...\\
\hspace*{3mm}CorrelationBackend="mex\_new", MaxHours=12);
\end{quote}
Examples show option syntax only; no production run was launched for this note.

\subsection*{Reuse saved posterior parameters}
Set \code{EstimateParameters=false}, \code{ParameterSource="posterior"},
and \code{SmoothStates=true}. The code loads the newest compatible parameter
snapshot and its original prior calibration, then smooths using the available
history. Appended observations can change past states without re-estimating
the fixed parameters.

\textbf{The existing \code{ParameterSmoothing="draws"} alternative.}
If estimation and smoothing use identical dates, the code can use the joint
estimation draws directly: these target the joint posterior once adequately
converged. For loaded draws or an extended history, the existing smoother
cycles through parameter draws between state sweeps. That is a heuristic;
ordinary common-target chain diagnostics are not reported for it. The new
prior option on page 3 instead holds each parameter set fixed for an entire
conditional chain, with its own warm-up.

\subsection*{Where the prior hyperparameters come from}
The priors are \textbf{empirical Bayes}: some data calibrate them before the
sampler starts. This remains true in the new prior option.

\begin{tabularx}{\linewidth}{@{}p{36mm}Y@{}}\toprule
Quantity & Calibration used by the code \\\midrule
$\bar B$, $\bar V_B$ & Fit an intercept-only OLS mean to complete returns in
the first 104 weeks. Set $\bar B$ to the sample mean and $\bar V_B=S_0/n_0$,
where $S_0$ is the covariance with 5\% diagonal shrinkage. \\
$\bar h$, $\bar r$ & From the same $S_0$: $\bar h=\log\operatorname{diag}(S_0)$;
standardize $S_0$ to a correlation matrix and take the lower off-diagonal
entries of its matrix logarithm for $\bar r$. \\
$V_{0B}$, $\nu_B$ & $\nu_B=n_0$ complete calibration returns and
$V_{0B}=\bar V_B\nu_B(0.01)^2$. \\
$a$, $b_h$, $b_r$ & Use 104-week rolling covariance windows through the
calibration cutoff. Compute changes in rolling log variances and matrix-log
correlations; take their time-series variances, then the median across
coordinates within each block. These set blockwise prior means $\mu_h,\mu_r$
(floored at $10^{-8}$). Set $a=10$ and $b=(a-1)\mu$. \\\bottomrule
\end{tabularx}
Here $n_0$ counts complete rows in the reserved window. Initial-state priors
are $B_1\sim\Normal(\bar B,4\bar V_B)$ independently of $V$ and
$h_0\mid\sigma_h^2\sim\Normal(\bar h,10\sigma_h^2)$,
$r_0\mid\sigma_r^2\sim\Normal(\bar r,10\sigma_r^2)$ componentwise.
Integrating the first random-walk step gives $11\sigma^2$ for the first
modeled $h$ or $r$ state. The rolling calibration includes later modeled
returns through the cutoff, so the exercise is not strictly data-free.

\clearpage
{\color{ink}\Large\bfseries Prior parameter draws + conditional smoothing\par}
\subsection*{3. The new option and its target}
Let $X=(B,h,r)$ and let $\pi_{\rm EB}(\theta)$ denote the calibrated parameter
prior. The new option targets
\[
\boxed{q(\theta,X\mid y)=\pi_{\rm EB}(\theta)\,p(X\mid y,\theta).}
\]
Returns enter the conditional smoother for $X$. They do not reweight
$\theta$ by its likelihood. Consequently $q$ is \textbf{not the joint
posterior} $p(\theta,X\mid y)$, and it is \textbf{not a prior-predictive
simulation}, which would generate new data instead of smoothing the observed
returns.

\begin{enumerate}
\item Draw $K=\code{PriorParameterDraws}$ independent parameter sets:
\[
V^{(k)}\sim\IW(V_{0B},\nu_B),\quad
\sigma_{h,j}^{2(k)}\sim\IG(a,b_h),\quad
\sigma_{r,l}^{2(k)}\sim\IG(a,b_r).
\]
The blocks and scalar variance draws are mutually independent.
\item For each $k$, freeze $\theta^{(k)}$ and start
$C=\code{NumChains}$ conditional state chains. Each chain warms up separately,
then retains $R=\code{RetainedDraws}$ states from
$p(X\mid y,\theta^{(k)})$ using the existing $B$, $r$ and $h$ updates.
\item Check diagnostics \emph{within each fixed parameter set}. Retain the
parameter-draw identifiers so conditional convergence and prior mixing
remain distinct.
\item With complete equal quotas, estimate a state summary by giving each
parameter set weight $1/K$:
\[
\widehat{E_q[f(X)]}=\frac1K\sum_{k=1}^{K}
\frac1{CR}\sum_{c=1}^{C}\sum_{s=1}^{R}f(X^{(k,c,s)}).
\]
\end{enumerate}
Do not redraw $\theta$ after one state sweep: a single sweep is generally not
a draw from the new conditional smoother.

\begin{quote}\small\ttfamily
result = run\_weekly\_model\_report( ...\\
\hspace*{3mm}ParameterSource="prior", PriorParameterDraws=10, ...\\
\hspace*{3mm}SmoothStates=true, WarmupIterations=100, ...\\
\hspace*{3mm}RetainedDraws=100, NumChains=4, ParallelChains=true, ...\\
\hspace*{3mm}CorrelationBackend="mex\_new", MaxHours=12);
\end{quote}
This requests \textbf{10 independent parameter sets} and
$10\times4\times100=4{,}000$ conditional state draws. All parameter sets
share one \code{MaxHours} stage budget. An incomplete run with unequal
component quotas is not pooled as an equally weighted mixture.

\subsection*{How to judge the result}
For each fixed $\theta^{(k)}$, the reference checks are
$\widehat R<1.01$, bulk ESS and tail ESS $\geq400$, and
MCSE/posterior SD $\leq0.05$. Enough retained draws and multiple independent
chains are needed for useful checks. Never compute cross-chain $\widehat R$
by treating different $\theta$ sets as chains of the same target. Good
conditional diagnostics do not establish accuracy of the outer prior
integral; $K$ controls a separate integration error.

\textbf{Scope of this version.} Prior mode skips posterior parameter
estimation and loading, requires \code{SmoothStates=true}, and does not
support automatic extension or resume. The existing publication builder
supports posterior runs with smoothing; this standalone note explains all
three modes. State figures and saved diagnostic artifacts remain available
for prior mode.

{\footnotesize\color{ink}Implementation basis: \code{run\_weekly\_model\_report.m},
\code{dsc\_prepare\_inference.m}, \code{dsc\_run\_inference.m},
\code{dsc\_sample.m}, and \code{dsc\_calibrate\_priors.m} in this branch.
No performance, posterior, or convergence results are claimed here.}
\end{document}
